\documentclass[12pt]{article}
\usepackage[T1]{fontenc}
\usepackage[utf8]{inputenc}
\usepackage{lmodern}
\usepackage{textcomp}
\usepackage{enumitem}
\usepackage{geometry}
\usepackage{graphicx}
\usepackage{amsmath, amssymb}
\geometry{margin=1in}
\begin{document}
\begin{center}
Chemistry - Undergraduate Level Assessment 7 \
Total Marks: 40 \
Answer all questions.
\end{center}

\section*{Multiple Choice (10 marks)}
\begin{enumerate}[label=\arabic*.]
\item A molecule has a rotational correlation time of 1 ns. At what magnetic field would protons in this molecule have the fastest spin-lattice relaxation? [Use eqns 5.3 and 5.4.]
\begin{enumerate}[label=(\alph*)]
    \item 3.74 T
    \item 5.19 T
    \item 6.08 T
    \item 9.49 T
\end{enumerate}
\item Infrared (IR) spectroscopy is useful for determining certain aspects of the structure of organic molecules because
\begin{enumerate}[label=(\alph*)]
    \item all molecular bonds absorb IR radiation
    \item IR peak intensities are related to molecular mass
    \item most organic functional groups absorb in a characteristic region of the IR spectrum
    \item each element absorbs at a characteristic wavelength
\end{enumerate}
\item Of the following compounds, which has the lowest melting point?
\begin{enumerate}[label=(\alph*)]
    \item HCl
    \item AgCl
    \item CaCl 2
    \item CCl 4
\end{enumerate}
\item Which nuclide has an NMR frequency of 115.5 MHz in a 20.0 T magnetic field?
\begin{enumerate}[label=(\alph*)]
    \item 17 O
    \item 19 F
    \item 29 Si
    \item 31 P
\end{enumerate}
\item Of the following ions, which has the smallest radius?
\begin{enumerate}[label=(\alph*)]
    \item K+
    \item Ca$_{2}$+
    \item Sc$_{3}$+
    \item Rb+
\end{enumerate}
\end{enumerate}

\section*{True / False (10 marks)}
\textit{Answer True or False}
\begin{enumerate}[label=\arabic*.]
\setcounter{enumi}{5}
\item True or False: The Q-factor for an X-band EPR cavity with a resonator bandwidth of 1.58 MHz is Q = 6012.
\item True or False: Tetramethylsilane is used as a reference because its protons are weakly shielded, which affects the chemical shift.
\item True or False: Doubling the volume of a buffer solution significantly increases its pH by diluting the weak acid and conjugate base concentrations.
\item True or False: Silicon is most commonly found in nature as a metallic element, occurring in its pure form in the earth's crust.
\item True or False: CH 3 Cl has a 1 H resonance approximately 1.55 kHz downfield from TMS in a 12.0 T magnetic field.
\end{enumerate}

\section*{Long Form Response (20 marks)}
\textit{Answer each question with a clear and complete explanation.}
\begin{enumerate}[label=\arabic*.]
\setcounter{enumi}{10}
\item Discuss the factors influencing the equilibrium polarization of an electron and a proton placed in a 1.0 T magnetic field at 300 K, and calculate the ratio of their polarizations.
\item Discuss the process by which the polarization of 13 C nuclei can be increased in a 600 MHz spectrometer. How long will it take for the polarization to reach twice its thermal equilibrium value at 298 K, given a T$_{1}$ of 5.0 s?
\end{enumerate}

\vfill
\noindent\textit{End of Paper}
\end{document}